\section{The second term and the singularity at one in indefinite signature}\label{sec:spinboundary}
Fix a base index $d$ as in Theorem~\ref{th:indmean}, a local-type
tuple at $d$, and its reference child $K$. Take $C=C(K)$ for proper
equivalence and $C=C(K)/\langle\delta_K\rangle$ for ordinary
equivalence. Let $T\subseteq C$ be the nonempty image of $T(K,L)$;
in general $T$ is a subset rather than a subgroup. Write these
groups additively. For $p\notin S_d$, let $v_p\in C$ be the class of the
idele equal to $p$ at $p$ and to $1$ elsewhere. Characters of $C$ will be denoted by $\chi$.

\begin{lemma}\label{lem:projection}
Give each $p\mid t$ one of the three local labels of Lemma~\ref{lem:indlocal}, and let $D\subset C$ be the subgroup generated by the $v_p$ whose label is isotropic. The number of actual $H$-classes in this extended tuple is $|T\bmod D|$.
\end{lemma}
\begin{proof}
At a nonisotropic label the full stabilizer of the child $U\perp p^2\langle u\rangle$ has spinor norms of even valuation: successively fix a unit diagonal basis of the unimodular $U$ by the unit reflections from Lemma~\ref{lem:indlocal}, then the unary factor also has norm of even valuation. The proper stabilizer contains all units. At an isotropic label its proper stabilizer has full spinor norm, so the new proper genus group is the old one modulo exactly the listed uniformizers. The other local stabilizers have not changed.

At each new prime the transporters of that local type are $SO(L_p)SO(K'_p)$ for the extended child $K'$. Their image is trivial. Thus the represented spinor labels are the quotient image of the old ones. The reversal translation descends to the same quotient: at every new prime choose an improper isometry preserving both the parent and the new child, by reflection in the unit summand perpendicular to the chosen nonisotropic line or to the chosen hyperbolic plane. Such choices can already be used to compute $\delta_K$. This proves the description for either group.
\end{proof}

For a subgroup $W\subset C$ put $\eta=[C:W]^{-1}$. A character of $C$ gives at the good primes a quadratic Dirichlet character $\chi(p)=\chi(v_p)$. These characters factor modulo $Q_d=8\prod_{p\in S_d,\,p>2}p$.
Dividing a uniformizer idele by the corresponding rational prime
leaves unit classes at $S_d$, determined modulo $8$ or $p$. Conversely unit tuples generate $C$ and all unit tuples are attained by prime progressions, so a nontrivial character gives a nonprincipal Dirichlet character. Characters are extended by zero at primes of the modulus. Define
\[
 J_W(s)=\prod_{p\notin S_d}\bigl(1+(2+\mathbf1_{v_p\in W})p^{-s}\bigr),\qquad
 k_W=\lim_{x\downarrow0}x^{2+\eta}J_W(1+x).
\]
The Euler product counts squarefree labelled extensions with
isotropic primes mapping into $W$. Character orthogonality yields
\[
 J_W(s)=\zeta(s)^{2+\eta}
       \prod_{\substack{\chi\in W^\perp\\\chi\ne1}} L(s,\chi)^\eta\,A_W(s).
\]
On the real axis $s>1$ the powers take their positive values. The Euler factors of $A_W$ are the quotients of the corresponding local factors. Outside $S_d$ their linear terms cancel, so the product is absolutely convergent near $1$, with positive value there. At the finitely many discarded primes the residual factor is $(1-p^{-s})^{2+\eta}$. Dirichlet's theorem $L(1,\chi)>0$ for nonprincipal quadratic characters now gives $k_W=A_W(1)\prod_{\chi\ne1}L(1,\chi)^\eta$, with characters restricted as in the display.

\begin{theorem}\label{th:indsecond}
For either orientation group in every indefinite rank $n\ge3$ there is a finite, strictly positive constant
\[
 K_H=\sum_d\frac1d\sum_{\text{base tuples at }d}
  \ \sum_{\substack{\chi\in\widehat C,\ \chi\ne1\\
                      \chi\ \text{nonconstant on }T}} k_{\ker\chi}.
\]
Writing $Z_{\mathrm{gen}}(s)=\sum_m c(m)m^{-s}=\zeta(s)^3V(s)$, one has
\begin{align*}
 Z_H(1+x)-Z_{\mathrm{gen}}(1+x)&\sim K_H x^{-5/2}\quad(x\downarrow0),\\
 \sum_{m\le X}a_H(m)&=\frac{C_L}2X(\log X)^2+
 \frac{K_H}{\Gamma(5/2)}X(\log X)^{3/2}+o_L(X(\log X)^{3/2}).
\end{align*}
Consequently no meromorphic germ at $1$ agrees with $Z_H$ on a right-hand neighbourhood.
\end{theorem}
\begin{proof}
\emph{Separating the first obstruction.}
For fixed $d$ and base tuple expand the function $f(D)=|T\bmod D|$ on the finite set of subgroups as $\sum_{D\subseteq W}\beta(W)$. The coefficients are determined descending from $\beta(C)=1$. For a hyperplane $W=\ker\chi$, $\beta(W)$ is $1$ if $\chi$ takes both signs on $T$, and zero otherwise. Thus the labelled-extension series is $\sum_W\beta(W)J_W(s)$. Subtracting its genus part removes the term $W=C$. Hyperplanes give exponent $5/2$ and the remaining proper subgroups give exponents at most $9/4$. This gives the real-axis asymptotic for each fixed tuple.

\emph{From Euler products to unsmoothed counts.}
For the per-tuple counting limit factor $J_W(s)=\zeta(s)^2 H_W(s)B_W(s)$ where
\[
 H_W(s)=\prod_{\substack{p\notin S_d\\v_p\in W}}(1-p^{-s})^{-1}.
\]
The Dirichlet coefficients of $H_W$ are nonnegative, and $B_W$ has absolutely convergent Dirichlet series on $\Re s>1/2$ by cancellation of the linear terms. At $1$ it is positive. Writing $h(m)$ for the coefficients of $H_W$, the real limit gives $H_W(1+x)\sim [k_W/B_W(1)]x^{-\eta}$. 

Apply Karamata's theorem for a Laplace--Stieltjes transform to the measure with mass $h(m)/m$ at $\log m$ \cite[Thm.~1.7.1]{BGT}. Its cumulative mass up to $\log y$ is asymptotic to $k_W(\log y)^\eta/(B_W(1)\Gamma(1+\eta))$. 

In the convolution with $d_2$, use $\sum_{j\le y}d_2(j)=y\log y+O(y)$ for $y\ge1$. The main sum divided by $y$ is the integral of that cumulative mass from $0$ to $\log y$, while the error divided by $y$ is bounded by the cumulative mass itself. Consequently after convolution with $B_W$ the summatory function of $J_W$ is asymptotic to $k_W y(\log y)^{1+\eta}/\Gamma(2+\eta)$. 
The last convolution follows by dominated convergence, using the resulting bound $O_W(y(1+\log y)^{1+\eta})$ before convolution and absolute convergence at $1$. The resulting estimate holds for each fixed tuple.

\emph{Uniformity over base indices.}
To sum over all base indices $d$, use the pointwise majorant
\[
 0\le |T\bmod D|-1\le |C/D|-1
 =\sum_{\chi\ne1}\mathbf1_{D\subseteq\ker\chi}.
\]
For each such hyperplane the coefficients of $J_{\ker\chi}$ are bounded by those of $\zeta(s)^2H_{\ker\chi}(s)$, by inspecting each nonnegative local factor. On real $1<\sigma\le2$ one has $H_{\ker\chi}(\sigma)^2\le\zeta(\sigma)L(\sigma,\chi)$, locally at split, inert and removed primes. Periodicity and zero mean give $L(\sigma,\chi)\ll1+\log Q_d$ uniformly: sum to $Q_d$ trivially and use partial summation with partial sums bounded by $Q_d$ for the tail. This proves $x^{5/2}J_{\ker\chi}(1+x)\ll\sqrt{1+\log Q_d}$. For summatory coefficients put $\sigma=1+1/\log y$ for $y\ge3$. Then $\sum_{m\le y}h(m)/m\le eH_{\ker\chi}(\sigma)$. The divisor convolution consequently bounds the count for a single hyperplane by $O(y(1+\log y)^{3/2}\sqrt{1+\log Q_d})$, enlarging the constant for small $y$. 
There are at most $4\cdot2^{|S_d|}$ characters. The additional square root of a logarithm is summable against $2^{|S_d|}c(d)/d$: the same Euler majorant converges with $d^{-(1-\epsilon)}$ for $0<\epsilon<1/2$. This justifies both limits by dominated convergence. For the genus term the error $O_L(X(1+\log X))$ follows from the divisor estimate in Theorem~\ref{th:indmean}, since $V$ has absolute convergence to the left of one.

In the construction of Theorem~\ref{th:indeuler}, the two classes in the same local tuple remain distinct even after ordinary orientation reversal. Hence its $T$ has at least two elements in $C$ for either choice of $H$. Some character of this finite binary vector space separates them. The corresponding hyperplane has positive coefficient, proving $K_H>0$. Finally $Z_{\mathrm{gen}}$ is meromorphic at one. If $Z_H$ were meromorphic there, their difference would have an integral order, incompatible with the positive $x^{-5/2}$ asymptotic.
\end{proof}
